%% sections/04-study.tex

\section{Study Design}
\label{sec:study}

\subsection{Corpus}

The corpus is the full Marathon D\&D Workshop: 49 sessions, 7 campaigns, 49 distinct
adventure modules. NPC arc-completions were logged systematically in session logs
from Campaign 3 onward (Campaigns 1--2 retroactively logged from session log text).

\subsection{Coding Protocol}

An arc-completion instance is logged when:

\begin{enumerate}
  \item An NPC delivers their designed arc-completion text (verbatim or close paraphrase).
  \item The session log records a party action immediately preceding the arc-completion
    that matches the firing condition.
  \item The DM does not provide additional NPC dialogue to explain or extend the moment.
\end{enumerate}

Instances that met criterion 1 but not criterion 2 (DM-scripted moments) were logged
separately as ``DM-initiated NPC moments'' and excluded from arc-completion counts.
Instances that met criterion 2 but not criterion 3 (extended DM narration after the
arc-completion) were logged as ``arc-completion with DM extension'' and counted as
partial instances.

In the full corpus, 40+ full instances and 7 partial instances were logged. No DM-initiated
NPC moments were logged (the AI system enforces firing-condition checking).

\subsection{NPC Classification}

NPCs were classified on two axes for cross-archetype analysis:

\begin{description}
  \item[Moral valence:] Protagonist-adjacent (ally, mentor, victim); Antagonist-adjacent
    (opponent, obstructor); Morally grey (professionally ambiguous, shifting allegiance).
  \item[Institutional role:] Grief-adjacent (NPC's arc explicitly tied to loss or
    named grief); Professional-code (NPC's arc tied to duty, craft, or institutional
    obligation); Situationally-triggered (NPC's arc tied to a specific situational
    condition not pre-figured in their backstory).
\end{description}

\subsection{Session Coverage Metric}

Session coverage is defined as the proportion of sessions in which at least one
full arc-completion fired. This metric captures the pattern's \textit{reliability}
rather than its frequency: a pattern that fires once in 88\% of sessions is more
practically useful than one that fires 5 times in 40\% of sessions.
